\section{Verwandte Arbeiten}
\label{sec:related}

Steiger et al. verbinden ein regelbasiertes Expertensystem mit mehreren
koordinierten Ansichten für die Überwachung eines Smart Grids
\cite{steiger2013smartgrid}. Eine Übersicht zeigt Zustände verteilter
Netzelemente; eine Analyseansicht macht den Verarbeitungspfad von Alarmen
nachvollziehbar. Gemeinsam mit diesem Projekt sind der Energiekontext,
Details-on-Demand und die Selektion eines Elements, die eine zweite Ansicht
aktualisiert. Der Unterschied liegt im Analyseziel: Steiger et al. unterstützen
operative Störungsdiagnose und erklären automatische Regeln. Die vorliegende
Anwendung arbeitet retrospektiv mit länderbezogenen Flüssen, Erzeugung und
Preis; sie enthält weder Alarmmodell noch Handlungsempfehlung.

Jäger et al. untersuchen die Kombination topologischer und geografischer
Information für das Management elektrischer Ausfälle
\cite{jaeger2016outage}. Ihre Designstudie zeigt, dass getrennte Räume einen
mentalen Abgleich verlangen, und nutzt interaktives Brushing-and-Linking. Das
stützt die Entscheidung, Länderbeziehung und Zeitmuster koordiniert zu zeigen.
Anders als dort ist Geografie für die Projektaufgabe keine zwingende
Datendimension: Die Abstände und Formen der Länder erklären weder Flussbetrag
noch Handel. Deshalb ersetzt ein kompakter radialer Graph die Kartenansicht.
Eine künftige Erweiterung könnte eine Karte ergänzen, wenn reale
Grenzkuppelleitungen und räumliche Engpässe einbezogen werden.

AMPLIO VQA von Stoffel et al. analysiert Simulationen für die Planung eines
Microgrid-Erzeugungsmixes \cite{stoffel2012amplio}. Beide Systeme behandeln
mehrdimensionale zeitbezogene Energiedaten und kombinieren Überblick,
Interaktion und Detail. AMPLIO vergleicht jedoch parametrisierte
Planungsszenarien und unterstützt visuelle Anfragen; dieses Projekt untersucht
gemessene historische Daten und einen nationalen grenzüberschreitenden
Austausch. Die gestapelte Zeitreihe ähnelt dem gemeinsamen Bedarf, einen Mix
über Zeit zu verstehen, während Netzwerk und Pixelmatrix speziell aus den
Dimensionen Partnerland und Stunde abgeleitet sind.

Methodisch steht die Anwendung in der Tradition koordinierter multipler
Ansichten. North und Shneiderman beschreiben Koordination als Grundlage für
Brushing-and-Linking und Overview-and-Detail \cite{north2000snap}. Die
vorliegende Implementierung erlaubt Nutzenden zwar nicht, Ansichten frei zu
konfigurieren, realisiert aber eine feste, domänenspezifische Kopplung über
Land und Zeit. Die Pixelmatrix folgt dem Prinzip pixelorientierter Techniken,
möglichst viele Einzelwerte gleichzeitig auf Bildschirmfläche abzubilden
\cite{keim2000pixel}; anders als allgemeine hochdimensionale Verfahren nutzt
sie eine fachlich feste Ordnung der Achsen.

Insgesamt unterscheidet sich das Projekt von den drei Energiearbeiten durch
die konkrete Kombination aus bilateralem Ländergraph, Erzeugungs-/Flusszeitreihe
und Länder-Stunden-Matrix auf realen Energy-Charts-Daten. Es ist kleiner und
nicht als Leitwartenwerkzeug evaluiert. Sein Beitrag liegt in einer
nachvollziehbaren, auf eine Lehr- und Explorationszielgruppe zugeschnittenen
Integration, nicht in einem neuen Visualisierungsalgorithmus.
